%\documentclass[11pt,a4paper,twoside]{article}

\usepackage[utf8]{inputenc}

\usepackage[T1,T2A]{fontenc}

\usepackage[english.ancient,russian.ancient]{babel}

\usepackage{graphicx}

\usepackage{array}

\usepackage[doi]{samgtu-bib}

\usepackage{samgtu}

\usepackage{samgtu-name}

\usepackage{mathtools}

\mathtoolsset{showonlyrefs}

\usepackage{desclist}

\usepackage{euscript}

\usepackage{mathrsfs} 

\usepackage{multicol}

\usepackage{mathrsfs}

\usepackage{cite}

\usepackage{cmap}

\usepackage{qrcode}

\allowdisplaybreaks[4]

\usepackage[nodayofweek]{datetime}

\usepackage[thinlines]{easybmat}



\renewcommand{\JournalYear}{2018}

\renewcommand{\JournalVolume}{22}

\renewcommand{\JournalNumber}{4}



\newdate{JournalDateSign}{29}{12}{2018}% Дата подписания Вестника в печать

\renewcommand{\JournalDateSign}{\displaydate{JournalDateSign}}

\newdate{JournalDatePrint}{16}{01}{2019}% Дата выхода Вестника из печати

\renewcommand{\JournalDatePrint}{\displaydate{JournalDatePrint}}

%\renewcommand{\JournalUPL}{16.56} % Усл. печ. л.

%\renewcommand{\JournalUIL}{16.53} % Уч.-изд. л.

\renewcommand{\JournalUPL}{15.24} % Усл. печ. л.

\renewcommand{\JournalUIL}{15.21} % Уч.-изд. л.

\renewcommand{\JournalRegNumber}{220/18} % Регистрационный номер

\renewcommand{\JournalOrderNumber}{2} % Номер заказа



%\renewcommand{\JournalRMonth}{Март} % Месяц выхода Вестника

%\renewcommand{\JournalEMonth}{March} % Месяц выхода Вестника

%\renewcommand{\JournalRMonth}{Июнь} % Месяц выхода Вестника

%\renewcommand{\JournalEMonth}{June} % Месяц выхода Вестника

%\renewcommand{\JournalRMonth}{Сентябрь} % Месяц выхода Вестника

%\renewcommand{\JournalEMonth}{September} % Месяц выхода Вестника

\renewcommand{\JournalRMonth}{Декабрь} % Месяц выхода Вестника

\renewcommand{\JournalEMonth}{December} % Месяц выхода Вестника



\newcommand{\totop}[1]{\raisebox{6pt}[1pt][2pt]{#1}}

\newcommand{\tottop}[1]{\raisebox{12pt}[1pt][2pt]{#1}} 

\newcommand{\totttop}[1]{\raisebox{18pt}[1pt][2pt]{#1}} 

\newcommand{\tottttop}[1]{\raisebox{24pt}[1pt][2pt]{#1}} 

\newcommand{\totttttop}[1]{\raisebox{30pt}[1pt][2pt]{#1}} 

\newcommand{\tobot}[1]{\raisebox{-6pt}[1pt][2pt]{#1}} 

\newcommand{\tobbot}[1]{\raisebox{-12pt}[1pt][2pt]{#1}} 

\newcommand{\tobbbot}[1]{\raisebox{-18pt}[1pt][2pt]{#1}}



\graphicspath{{./00PIC/}}

%

\vsgtu{1614} %registration number

\FirstPage{430}

\LastPage{446}

%

%\pdfcrop

%\draft

%

\ResearchArticle

%

%\begin{document}





\hfill \qrcode[hyperlink,height=.8in]{http://doi.org/10.14498/vsgtu1614}

\vspace{-.8in}



\udc{539.376}

\msc{74S30, 74D10}



\rutitle{Численное и экспериментальное исследование чистого~изгиба балок из титанового сплава АБВТ-20~в~условиях ползучести с~учетом различных~свойств на растяжение и~сжатие}

\rutitleshort{Численное и экспериментальное исследование чистого изгиба балок из титанового сплава\dots}



\entitle{Numerical and experimental research of pure bending of~beams  made of the titanium ABVT-20 alloy  with different properties for tension and compression under~creep conditions}

\entitleshort{Numerical and experimental research of pure bending of beams\dots}



\ruauthor{С.~В.~Иявойнен, А.~Ю.~Ларичкин, В.~Е.~Колодезев}{И\,я\,в\,о\,й\,н\,е\,н\,~С.~В.,  Л\,а\,р\,и\,ч\,к\,и\,н\,~А.~Ю., К\,о\,л\,о\,д\,е\,з\,е\,в\,~В.~Е.}

\enauthor{S.~V.~Iyavoynen, A.~Yu.~Larichkin, V.~E.~Kolodezev}{I\,y\,a\,v\,o\,y\,n\,e\,n~S.~V., L\,a\,r\,i\,c\,h\,k\,i\,n~A.~Yu., K\,o\,l\,o\,d\,e\,z\,e\,v~V.~E.}



\ruaffil{\ruaffid{982}{Институт гидродинамики им. М.А. Лаврентьева СО РАН\\ 

Россия, 630090, Новосибирск, Лаврентьева проспект, 15}.

%\and 

%Новосибирский государственный университет,\\ 

%Россия, 630090, Новосибирск, ул. Пирогова, 2.

}



\enaffil{\enaffid{982}{Lavrentyev Institute of Hydrodynamics \\ 

of Siberian Branch  of the Russian Academy of Sciences, \\ 

15, Lavrentyev av., Novosibirsk, 630090, Russian Federation}.

%\and 

%Novosibirsk State University, 2, Pirogova st., Novosibirsk, 630090, Russian Federation.

}



\ruauthorsdetails{3}{% Количество авторов

\fullauthor{\rupid{135788}{Светлана Владимировна Иявойнен}}

\CAuthor % Автор-корреспондент

\orcid{0000-0002-1478-0533} % ORCID ID автора 

\regalia{аспирант} 

\position{младший научный сотрудник} 

\dept{лаб. статической прочности} 

\email[b]{svetaiyavoynen@gmail.com}

\fullauthor{\rupid{43295}{Алексей Юрьевич Ларичкин}}

\orcid{0000-0002-7306-9522} % ORCID ID автора 

\regalia{кандидат физико-математических наук} 

\position{научный сотрудник} 

\dept{лаб. статической прочности} 

\email{larichking@gmail.com}

\fullauthor{\rupid{143089}{Вадим Евгеньевич Колодезев}}

\orcid{0000-0001-8035-7453} % ORCID ID автора 

\regalia{кандидат технических наук} 

\position{инженер-технолог} 

\dept{лаб. статической прочности} 

\email[b]{kolodezev.vadim@yandex.ru}

}



\enauthorsdetails{3}{

\fullauthor{\enpid{135788}{Svetlana V. Iyavoynen}}

\CAuthor % Сorresponding Author

\orcid{0000-0002-1478-0533} % ORCID ID

\regalia{Postgraduate Student} 

\position{Junior Researcher} 

\dept{Lab. of Static Strength}

\email[b]{svetaiyavoynen@gmail.com}

\fullauthor{\enpid{43295}{Aleksey Yu. Larichkin}}

\orcid{0000-0002-7306-9522} % ORCID ID avtora 

\regalia{Cand. Phys. \& Math. Sci.} 

\position{Researcher} 

\dept{Lab. of Static Strength}

\email[br]{larichking@gmail.com}

\fullauthor{\enpid{143089}{Vadim E. Kolodezev}}

\orcid{0000-0001-8035-7453} % ORCID ID avtora 

\regalia{Cand. Tech. Sci.} 

\position{Incorporated Engineer} 

\dept{Lab. of Static Strength}

\email[br]{kolodezev.vadim@yandex.ru}

}



\rukeywords{высокотемпературная ползучесть, разносопротивляемость, чистый изгиб, экспериментальные исследования, моделирование, сплав АБВТ-20}

\enkeywords{high temperature creep,

different resistivity,

pure bending,

experimental studies,

modeling,

ABVT-20 alloy}





\ruabstract{Рассматривается решение задачи чистого изгиба балки прямоугольного сечения в режиме ползучести с учетом различных свойств ползучести на растяжение и сжатие. Построен алгоритм и разработано программное обеспечение для математического моделирования процесса перераспределения напряжений по высоте балки с~учетом накопления повреждений. Моделирование процессов ползучести разупрочняющегося материала происходит на основе уравнений кинетической теории ползучести и повреждаемости. Численное решение задачи проводится на основе метода Рунге"--~Кутты"--~Мерсона. Проведено сравнение результатов моделирования с экспериментальными данными чистого изгиба балок прямоугольного сечения из титанового сплава АБВТ-20 при действии знакопеременного изгибающего момента в условиях продолжительного воздействия  температуры ($750^\circ$C), которое показало удовлетворительное соответствие результатов расчета экспериментальным данным.}



\enabstract{Solution of the problem of pure bending of a beam of rectangular cross section taking into account the difference of properties tension and compression under creep is considered. Program algorithm of mathematical simulation of the stress redistribution process along the height of a beam with allowance for damage accumulation is constructed and implemented. Modeling of creep processes of softening material is based on equations of the kinetic theory of creep and damage. In this paper, Runge–Kutta–Merson numerical integration algorithm for creep damage analysis is presented. The simulation results are compared with the experimental data of pure bending of rectangular section beams from the titanium ABVT-20 alloy under the action of an alternating moment and a prolonged exposure to temperature of 750$^\circ$C. A~satisfactory agreement between the simulation results and the experimental data was obtained, taking into account the duration of the temperature aging in the creep law.}



\newdate{iyava}{06}{03}{2018}

\date{\displaydate{iyava}}

\newdate{iyavb}{24}{08}{2018}

\revisiondate{\displaydate{iyavb}}

\newdate{iyavc}{03}{09}{2018}

\accepteddate{\displaydate{iyavc}}



\newdate{iyave}{10}{10}{2018}

\renewcommand{\JournalDataOnline}{\displaydate{iyave}}



%\ForPeerReview

\makerutitle





\Section[n]{Введение}

Сохранение ресурса и свойств материала изделия на стадии его производства является актуальной проблемой современного авиастроения. Одним из возможных путей решения является внедрение в производство технологий формообразования, основанных на явлении ползучести и сверхпластического течения материала. Такие технологические процессы широко внедряются в производство зарубежными авиапроизводителями (Airbus, Boeing)~\cite{iya:Ribeiro,iya:Beal}, поскольку позволяют сокращать количество технологических циклов обработки изделия на стадии формообразования, совмещать процессы медленного деформирования и старения, сохранять ресурс пластической деформации. Широкое распространение технологий медленного горячего формообразования затруднено наличием малого количества исследований о влиянии условий технологических процессов на прочностные свойства материалов в готовом изделии (прочность, долговечность, ударная вязкость, трещиностойкость и т.~п.). Отсутствие информации об оптимальных режимах деформирования различных конструкционных материалов также не способствует широкому внедрению таких технологий.

Варианты технологии горячего формообразования в режимах ползучести применяются на нескольких авиастроительных заводах России.



В США и Великобритании процесс формовки при температуре старения начали применять в восьмидесятых годах XX века. 

В~иностранной литературе рассматриваемый процесс называют ``creep age forming'' (сокращенно CAF)~\cite{iya:Ribeiro,iya:Beal, iya:add:1, iya:add:2, iya:add:3 }. 

Целью применения технологии являлась формовка больших алюминиевых панелей (около 15~м) для получения профиля крыла и поверхностей сложной геометрии. Технологию формовки применяли к элементам конструкций самолетов Airbus 330 и 340, придавая кривизну стрингерам при помощи кручения и изгиба~\cite{iya:Ribeiro}. 

В~работе~\cite{iya:Ribeiro} приведен обзор ключевых зарубежных компаний, развивающих подход формовки в~режимах старения, а~также представлен один из подходов к~формообразованию деталей в~этом режиме при помощи вакуумизации, когда для прижатия заготовки к~оснастке используют вакуумный мешок"=диафрагму. 

Авторы подчеркивают, что после формовки прямоугольной пластины в~цилиндрическую поверхность деталь испытывает около 70\% распружинивания. 

Показано, что при растяжении образцов из алюминиевого сплава 7075 (аналог отечественного сплава В95) в~режиме старения остаточные напряжения не превосходят 25~МПа по сравнению с~напряжениями в~150~МПа при холодной вытяжке. Авторы предлагают для создания больших оребренных панелей сложной геометрии отдельно формовать ребра и~саму панель, а~затем сваривать их между собой. Чтобы избежать искажений поверхности после распружинивания, используются различные методы, например термофиксация.



В работе~\cite{iya:Oleinikov} приводится решение обратной задачи формообразования крупногабаритных монолитных разнотолщинных панелей двойной кривизны из сплава В95очТФ (состояние поставки) в~режиме термофиксации для температуры старения, равной 165\textcelsius.  

Аббревиатура <<оч>> означает <<очень чистый>>, которая указывает на низкое содержание кремния в сплаве (до 0.1\%, что в~пять раз меньше, чем в~сплаве В95). 

Стоит отметить, что в~\cite{iya:Oleinikov} учтено изменение механических свойств сплава В95очТФ в~процессе термической фиксации. 

Неучет такой эволюции свойств завышает максимальные напряжения в~формуемой панели почти в~три раза. 

Приводятся вариационная формулировка задачи и~результаты конечноэлементного моделирования формообразования панели крыла в среде {\tt MSC.Marc}. 

Решения, основанные на выводах из данной работы, применяются для серийного формообразования крыловых панелей с~использованием эффектов ползучести и~деформационного старения на <<КнААЗ>> филиала ПАО <<Компания <<Сухой>>. 



Алгоритмы расчета формообразования деталей усложняются в случае модели материала, разносопротивляющегося растяжению и сжатию при ползучести.

В~\cite{iya:Korobeynikov} решены трехмерные задачи по кручению металлических пластин в условиях ползучести под действием постоянных сосредоточенных сил, приложенных в ее углах. Представлен алгоритм определения компонент тензора напряжений, реализованный в модели материала конечноэлементного пакета {\tt PIONER} (разработка ИГиЛ СО РАН) для определяющих соотношений ползучести с~учетом разных свойств материала при растяжении и~сжатии. 

Приведено сравнение результатов трех случаев моделирования кручения толстой плиты с~экспериментальными данными. 

В~первом случае в~законе ползучести использовались параметры, полученные из одноосных экспериментов только на сжатие, во втором "--- только на растяжение, в~третьем "---  обе группы параметров. 

Использование модели с~учетом различия свойств на растяжение и~сжатие по сравнению с моделью, где учитываются свойства ползучести только на растяжение или только на сжатие, позволяет добиться удовлетворительного соответствия расчетов и~данных эксперимента, что увеличивает точность формы упреждающей оснастки. 

Отметим, что модель разносопротивляемости внедрена в~конечноэлементный пакет {\tt PIONER}, что позволяет решать трехмерные задачи ползучести.



Монография~\cite{iya:Lok2016} посвящена фундаментальному описанию явления ползучести. 

В~ней приведены решения различных задач длительной прочности с~приложением к~технике, в~том числе задач изгиба балок, пластин, оболочек. 

Представлены подходы к~описанию ползучести и~накоплению повреждений при сложном напряженном состоянии. 

В~частности, получены решения задачи чистого изгиба балок в~режимах ползучести с~учетом различных свойств на растяжение и~сжатие, а~также поврежденности материала. Зависимость скорости ползучести и скорости изменения сплошности от напряжений принимается в виде дробно"=степенных функций. Приводится решение с~учетом фронта разрушения.



В~\cite{iya:Beal} приводятся различные способы формовки титановых сплавов, в~том числе в~режимах ползучести и~сверхпластичности. 

В~работе~\cite{iya:Beal} отмечено, что использование явления сверхпластичности для формообразования деталей истребителя F-15 приводит к~снижению их стоимости на 58\% и~снижению массы на 31\%. 

Приводится пример, что ранее деталь двигателя гондолы для самолета Boeing 757 изготовливалась из сорок одной детали и~более двухсот крепежных узлов (материал Ti-6Al-4V), а~с~применением формовки в~режиме сверхпластического течения эта деталь формообразуется из одного листа.



Однако добиться режимов сверхпластического течения не так просто: наравне с~основными параметрами процесса формовки "--- температурой и скоростью деформирования "--- рядом авторов отмечается важность влияния размера зерна материала~\cite{iya:Kaibishev}. 

Управление этими тремя параметрами позволяет деформировать материалы в~режимах сверхпластичности. 

Управление предполагает подготовку зерна материала, знание температуры динамической рекристаллизации и~регулирование скорости формовки. Размер зерна существенно влияет на механические свойства материала. Уменьшение зерна титанового сплава ВТ-6 (Ti-6Al-4V) до 10~мкм после всесторонней ковки дает понижение температуры сверхпластического течения и~усиление эффекта диффузионной сварки~\cite{iya:Safiullin}. 

На основе этих научных достижений удалось создать технологию получения полой лопатки турбины для вентилятора двигателя ПД-14.



Эффекты, связанные с~процессом ползучести, используются и~для последующей обработки изделий. Упрочнение поверхностных слоев лопаток и~иных деталей двигателя проводится при помощи дробеструйной обработки, алмазного выглаживания, прокатки роликом, что увеличивает срок службы изделия. 

В~технике известно благоприятное влияние сжимающих напряжений в~поверхностных слоях деталей на их усталостную долговечность. 

Вопросы, связанные с~релаксацией напряжений в поверхностно"=упрочненных слоях элементов конструкций при ползучести, достаточно полно освещены в работе~\cite{iya:Radchenko}. Авторы приводят способы учета пластических деформаций на фоне процесса ползучести. Дается обширный материал описания кинетики процесса накопления повреждений при ползучести. Приводится метод восстановления картины напряжений в поверхностном слое детали цилиндрической формы и описание процесса релаксации напряжений по его глубине.



Эффекты, связанные с формообразованием в режиме ползучести, часто зависят от времени выдержки при температуре. В~настоящей работе рассмотрены особенности знакопеременного изгиба прямоугольных балок из титанового сплава АБВТ-20 в~режиме ползучести с~учетом температурной выдержки без нагрузки. 

В~данной работе используется степенной закон ползучести с~учетом поврежденности материала, 

в~отличие от работы~\cite{iya:Loko2012}, где используется дробно"=степенная зависимость скорости деформаций ползучести от напряжения. Кроме этого, в~настоящей работе в~реологической модели для исследуемого материала используется зависимость скорости ползучести от времени температурной выдержки, которая наблюдалась в~экспериментальных исследованиях~\cite{iya:Larichkin}.





\smallskip



\Section{Математическое моделирование чистого изгиба балок}

Рассмотрим процесс ползучести чистого изгиба прямоугольной балки шириной $b$ и высотой $h$. Изгибающий момент $M$ действует в~плоскости симметрии балки и~не приводит к~появлению пластических деформаций. 

Моделирование процесса проводится для материала, имеющего различные свойства ползучести на растяжение и~сжатие. Расчет основан на уравнениях энергетического варианта теории ползучести~[\citen{iya:Sosnin}--\citen{iya:Sosnin1970}].



Считаем, что полная деформация $\varepsilon^{total}$ в любой момент времени есть аддитивная составляющая упругой деформации $\varepsilon^{e}$ и~деформации ползучести~$\varepsilon$:

\begin{equation*}

\varepsilon^{total}=\varepsilon^{e}+\varepsilon=\frac{\sigma}{E} + \varepsilon,

		\end{equation*}

 где $\sigma(t)=P/S_0$  "--- номинальное действующее напряжение в~балке в момент времени $t$; 

 $P$ "--- нагрузка; $S_0$ "--- начальная площадь; $E$ "--- модуль упругости материала.

 

Для случая чистого изгиба справедлива гипотеза Бернулли:

\begin{equation}\label{iya:ursigma}

\varepsilon^{total}=\varkappa(t)(z+\delta(t)),\quad \varepsilon(0)=0,

		\end{equation}     

где $\varkappa(t)$ и $\delta(t)$ "--- кривизна балки и смещение нейтральной поверхности  для некоторого момента времени $t$ соответственно; $z$ "--- координата по высоте балки с~началом в~срединной поверхности.



Процесс ползучести материала c учетом его повреждаемости описывается уравнениями вида~[\citen{iya:Loko2012}--\citen{iya:Gorev1973}]:

\begin{equation}

\label{iya:sist}

\begin{array}{l}

\displaystyle 

\frac{dA}{dt}=

    \begin{cases}

    \displaystyle 

	\frac{B_{A_{1}}\sigma^{n_1}}{ (1-\omega)^{m_1}}, & \sigma>0,    \\[2mm]

    \displaystyle 	

	\frac{B_{A_{2}}|\sigma|^{n_2}}{ (1-\omega)^{m_2}},& \sigma\leq 0;

	\end{cases}

\\[5mm]

\displaystyle

	\frac{d\omega}{dt}=

    \begin{cases}

    \displaystyle

	\frac{B_{\omega_{1}}\sigma^{r_1}}{(1-\omega)^{m_1}},& \sigma>0,    \\[2mm]

	\displaystyle

	\frac{B_{\omega_{2}}|\sigma|^{r_2}}{(1-\omega)^{m_2}}, & \sigma\leq 0;

	\end{cases}

\\[5mm]

\displaystyle

	\frac{d\varepsilon}{dt}=\frac{dA}{dt}\frac{1}{\sigma},

     \end{array}

  \end{equation}

где $A$ "--- величина удельной работы рассеяния, $A(z_j,0)=0$; 

$\omega$ "--- параметр поврежденности, $\omega(z_j,0)=0$;  

$B_{A_{1}}$, $B_{\omega_{1}}$,  $n_{1}$,  $r_{1}$,  $m_{1}$ и $B_{A_{2}}$, 

$B_{\omega_{2}}$,  $n_{2}$,  $r_{2}$,  $m_{2}$ "--- положительные константы материала на растяжение (с~индексом~1) и сжатие (с~индексом~2).







Уравнения равновесия  имеют следующий вид:

	\begin{equation}\label{eq:ulg:urrav}

		b\int_{-h/2}^{h/2}\sigma (z+\delta) dz=M, \quad  \int_{-h/2}^{h/2}\sigma  dz=0.

		\end{equation}

Кривизна балки $\varkappa$ в любой момент времени находится из~\eqref{iya:ursigma},~\eqref{eq:ulg:urrav}:

	\begin{equation}\label{iya:kappa}

		\varkappa=\frac{M}{E J_{\delta}}+\frac{b}{ J_{\delta}}\int_{-h/2}^{h/2} \varepsilon (z+\delta)  dz,

		\end{equation}

где $J_{\delta}=b(h^3/12+h  \delta^2)$ "--- осевой момент инерции сечения балки с~учетом смещения нейтральной оси $\delta$.

        

Из условия отсутствия нормальных внешних усилий в сечении бруса определяется $\delta$:

	\begin{equation}\label{iya:urdelta}

		\delta=\frac{1}{\varkappa h}\int_{-h/2}^{h/2}\varepsilon  dz.

		\end{equation}

Подставляя значения напряжений в каждой точке разбиения по высоте поперечного сечения балки из~\eqref{iya:ursigma} в~\eqref{iya:sist}, получаем систему обыкновенных дифференциальных уравнений   относительно  удельной работы рассеяния $A$, повреждаемости $\omega$ и деформации $\varepsilon$. 



Численное моделирование изгиба разносопротивляющейся балки проводится интегрированием уравнений~\eqref{iya:sist} с использованием~\eqref{eq:ulg:urrav}–\eqref{iya:urdelta} посредством метода Рунге"--~Кутты"--~Мерсона~\cite{iya:Runge}. Метод основан на схеме Рунге"--~Кутты первого порядка итерационного решения задачи Коши, где определяются четыре коэффициента для нахождения неизвестных функций правых частей решаемой системы  уравнений. Схема Мерсона добавляет дополнительный пятый коэффициент, что позволяет определять погрешность решения на каждом шаге интегрирования по времени. 



\smallskip



\Section{Численный анализ чистого изгиба балки в условиях ползучести} 

 Разбивая сечение балки по высоте на $k$ равных интервалов, из~\eqref{iya:sist} получаем систему дифференциально"=алгебраических уравнений в каждой точке разбиения по высоте балки:

 \begin{equation}\label{iya:eq6}

   \begin{array}{l}

\displaystyle 

\frac{dA_{j}}{dt}=

    \begin{cases}

    \displaystyle 

	\frac{B_{A_{1}}\sigma_{j}^{n_1}}{ (1-\omega_{j})^{m_1}}, & \sigma_{j}>0, 

    \\[2mm]

    \displaystyle 

	\frac{B_{A_{2}}|\sigma_{j}|^{n_2}}{ (1-\omega_{j})^{m_2}},& \sigma_{j}\leq 0;

	\end{cases}

   \\[4mm]

\displaystyle 

	\frac{d\omega_{j}}{dt}=

    \begin{cases}

\displaystyle 

	\frac{B_{\omega_{1}}\sigma_{j}^{r_1}}{(1-\omega_{j})^{m_1}}, & \sigma_{j}>0, 

    \\[2mm]

\displaystyle 

	\frac{B_{\omega_{2}}|\sigma_{j}|^{r_2}}{(1-\omega_{j})^{m_2}},& \sigma_{j}\leq 0;

	\end{cases}

    \\

\displaystyle 

	\frac{d\varepsilon_{j}}{dt}=\frac{dA_{j}}{dt}\frac{1}{\sigma_{j}}, \quad  j=0,1,\dots,k.

     \end{array}

  \end{equation}

 

В начальный момент времени в каждом волокне сечения балки имеем

\begin{equation*}

	t=0:\	\varepsilon(z_{j},0)=0; \ \omega(z_{j},0)=0; \ A(z_{j},0)=0.  

		\end{equation*}

  Заменяя интегралы  в уравнениях~\eqref{iya:kappa} и \eqref{iya:urdelta} на конечные суммы по формуле Симпсона~\cite{iya:Simpson}, имеем    

\begin{equation}

\begin{array}{l}

\displaystyle

\varkappa(t)=\frac{M}{E  J_{\delta}} +\frac{b  h/k}{3  J_{\delta}}\sum\limits_{j=1}^{k-1} 

\bigl[\varepsilon(z_{j-1},t)z_{j-1}+ 4  \varepsilon(z_{j},t)z_{j}+\varepsilon(z_{j+1},t)z_{j+1}\bigr],

\\

\displaystyle

\delta(t)=- \frac{ h/k}{3 \varkappa h}\sum\limits_{j=1}^{k-1}

\bigl[\varepsilon(z_{j-1},t)z_{j-1}+ 4  \varepsilon(z_{j},t)z_{j}+\varepsilon(z_{j+1},t)z_{j+1}\bigr],

\end{array}

\label{iya:eq7}

\end{equation}

 где $z_{j}$ "--- значение высоты балки в $j$-той точке разбиения.

 

Для решения задачи ползучести балки используется широко применяемый метод <<шагами по времени>>. 

Для этого кроме дискретизации по пространственной координате $z$ задается дискретизация по времени, т.е. $t_{i+1}\hm =t_i+\Delta t_i$ ($t_0=0$), 

где $\Delta t_i$ "--- шаг (в общем случае "--- переменный) интегрирования системы дифференциальных уравнений \eqref{iya:eq6}, при этом начальное напряженное состояние на отрезке $t\in[t_i,t_{i+1}]$ привязывается к~времени $t=t_i$. 

Укрупненно алгоритм решения следующий: 

\begin{itemize}

\item[1)] задаются начальные данные $\varepsilon(z_j,0)=0$, $A(z_j,0)=0$, $\omega(z_j,0)=0$ и~величина момента $M$;

\item[2)] решается упругая задача чистого изгиба балки; поскольку разносопротивляемость заложена лишь для деформации ползучести, то вследствие симметрии упругих свойств на растяжение--сжатие начальная величина смещения нейтральной оси $\delta(0)=0$; 

\item[3)] с использованием распределения напряжений $\sigma(z_j,0)$ как начального из решения упругой задачи в соотношениях \eqref{iya:eq6} решается система дифференциальных уравнений \eqref{iya:eq6} методом Рунге"--~Кутты"--~Мерсона~\cite{iya:Runge}, находятся деформации в момент времени $t=t_{i+1}$ при $i=0$, 

далее из \eqref{iya:eq7} последовательно определяются $\varkappa(t_{i+1})$ и $\delta(t_{i+1})$ и, наконец,

$$\sigma(z_j,t_{i+1})=E \bigl[\varkappa(t_{i+1})\bigl(z_j+\delta(t_{i+1})\bigr)-\varepsilon(t_{i+1})\bigr].$$

Далее алгоритм повторяется при $t=t_{i+1}$, $i=2,3,\dots$.

\end{itemize}

 

 При реализации численного алгоритма использовалась заложенная в~\cite{iya:Runge} возможность автоматического выбора шага интегрирования $\Delta t_i$ (его увеличения и уменьшения в зависимости от требуемой погрешности вычислений).

 

\smallskip 

 





 \Section{Испытания на ползучесть сплава АБВТ-20 при температуре  750\,\textcelsius} 

На испытательных машинах лаборатории статической прочности ИГиЛ СО РАН проведена серия испытаний по растяжению и~сжатию круглых цилиндрических образцов из  сплава АБВТ-20 при температуре  750\,\textcelsius. 

Испытательная машина представляет собой установку рычажного типа с~максимальным усилием 50\,000~Н. 

Нагружение образцов, соединенных с коротким концом рычага и станиной машины, проводится за счет установки грузов на длинном конце рычага. 



Установки  укомплектованы нагревательными системами, которые состоят из регулятора нагрева и печного пространства. Нагрев производится при помощи галогеновых ламп, регулировка температуры проводится контрольно"=измерительным прибором ОВЕН посредством термопар типа хромель"=копель. 

Нагревательная система позволяет проводить испытания на ползучесть до 1\,000\,\textcelsius, после выхода на целевую температуру возможное отклонение не превышает $\pm3$\,\textcelsius. 



Измерение удлинений (укорочений) образцов проводилось посредством систем замеров с использованием электрических датчиков линейного перемещения, данные с которых записывались в память компьютера в виде таблиц. В силу длительности испытаний частота измерений устанавливалась один раз в три секунды.



Образцы для испытаний вырезались из плиты сплава АБВТ-20 толщиной $18.5$ мм. Диаметр испытываемых образцов составлял 8 мм, длина рабочей части $l_{0}$:  45 мм на растяжение, 16 мм на сжатие.   

Образцы разделялись на два типа: 

\begin{itemize}

\item[--] первый тип "--- образцы, вырезанные из плиты вдоль направления проката; 

\item[--] второй тип "--- образцы, вырезанные в поперечном направлении относительно проката.

\end{itemize} Образцы для испытаний изготавливались в ИГиЛ СО РАН.



При одноосном испытании к образцу прикладывалось постоянное напряжение $\sigma = P/S ={\rm const}$. Значение усилий корректировалось через каждые 0.5\,\% деформации из условия несжимаемости материала при неупругом деформировании. Несжимаемость предполагает неизменность объема образца $V_{0}=l_{0}S_{0}= l S$, отсюда, зная его удлинение $\Delta l$ в~процессе испытания, можно определять текущее значение площади сечения $S=l_{0}S_{0}/l=l_{0}S_{0}/(l_{0}+\Delta l)$ и~поддерживать необходимое усилие $P$. 

Стоит отметить, что принимая условие несжимаемости материала, мы считаем, что деформация ползучести равномерно распределена по длине образца, но в~одноосных испытаниях это не так. 

Почти всегда при растяжении, если, конечно, мы не попали в температурно"=скоростной режим сверхпластического деформирования, будет происходить локализация деформаций. 

По окончании установившейся стадии ползучести при растяжении будет возникать шейка. 

Кроме того, эффективная площадь поперечного сечения в~процессе растяжения не равна геометрической в силу накопленных повреждений. 

Ряд авторов указывает на противоречивость принятия гипотезы о~несжимаемости при оценке поврежденности материала и~предлагают использовать закон сохранения массы $\rho_{0}l_{0}S_{0}=\rho l S$, где $\rho_{0}$ и~$\rho$ "--- начальная и текущая плотности материала. 

Закон сохранения массы переходит в~условие несжимаемости при равенстве плотностей материала до и~пос\-ле деформирования $\rho_{0}=\rho$, что наблюдается при отсутствии разрыхления материала. 

Однако в~настоящей работе, следуя идеям Ю.~Н.~Работнова~\cite{iya:Rabotnov} и~О.~В.~Соснина~\cite{iya:Sosnin,iya:Gorev2014}, для корректировки усилия в~процессе деформирования мы использовали условие несжимаемости материала. 



На рис.~\ref{iya-1.eps} приведены зависимости деформации от времени для одноосных испытаний круглых образцов из сплава АБВТ-20 при температуре 750\,\textcelsius, экспериментальные данные отображены маркерами в координатах <<деформация~$\varepsilon$~-- время~$t$>>, где $\varepsilon=\ln(1+\Delta l/l_0)$. 

Стоит отметить, что для данной температуры материал проявляет слабые свойства анизотропии при ползучести, что отражается в~различии скоростей деформации при одинаковых напряжениях и в характере разрушения образцов при растяжении (см. рис.~\ref{iya-1.eps}, {\sl a--b}). Здесь символом $(*)$ обозначен момент разрушения образца "--- разделение его на две части. Также стоит отметит, что в испытаниях на растяжение при уменьшении прикладываемых напряжений имеет место рост предельных деформаций для образцов второго типа и уменьшение предельных деформаций для образцов первого типа. На рис.~\ref{iya-1.eps}, {\sl c--d} приведены экспериментальные данные испытаний на сжатие.





 \begin{figure}[b!]

 

\vspace{-5mm}



\footnotesize



\centerline{

\begin{tabular}{ccc}

\includegraphics[scale=1.6]{iya-1a} & &

\includegraphics[scale=1.6]{iya-1b} \\

\sl   a & & \sl   b \\

& & \\

\includegraphics[scale=1.6]{iya-1c} & &

\includegraphics[scale=1.6]{iya-1d} \\

\sl  c & & \sl   d \\

\end{tabular} 

}

% 

%\noindent

%\begin{minipage}[h]{0.48\textwidth}

%\centering  \scriptsize \sl

%\includegraphics[width=0.98\textwidth]{iya-1a}\\

%{\sl a}

%\end{minipage}

%\begin{minipage}[h]{0.48\textwidth}

%\centering  \scriptsize \sl

%\includegraphics[width=0.98\textwidth]{iya-1b}\\

%{\sl b}

%\end{minipage}

%

%\vspace{2mm}

%

%\noindent

%\begin{minipage}[h]{0.48\textwidth}

%\centering  \scriptsize \sl

%\includegraphics[width=0.98\textwidth]{iya-1c}\\

%{\sl c}

%\end{minipage}

%\begin{minipage}[h]{0.48\textwidth}

%\centering  \scriptsize \sl

%\includegraphics[width=0.98\textwidth]{iya-1d}\\

%{\sl d}

%\end{minipage}

\caption{Данные одноосных испытаний круглых образцов из сплава АБВТ-20 при температуре 750\,\textcelsius: ({\sl a}\/) растяжение образцов второго типа; ({\sl b}\/) растяжение образцов первого типа;

 ({\sl c}\/)~сжатие образцов второго типа; ({\sl d}\/) сжатие образцов первого типа

\label{iya-1.eps}}



\smallskip 

[Figure~\ref{iya-1.eps}.

The data for an uniaxial tests of the circle specimens at 750\,\textcelsius: ({\sl a}\/) tension specimens cut out crosswise of the rolling direction of plate;

({\sl b}\/)~tension  specimens cut out along the direction of the rolling direction of plate;

 ({\sl c}\/)~compression specimens cut out crosswise of the rolling direction of plate;

 ({\sl d}\/)~compression specimens cut out along the direction of the rolling \centerline{direction of plate]}

 

 \vspace{-5mm}

 

\end{figure}



\begin{figure}[h!]

\noindent

\begin{minipage}[h]{0.56\textwidth}



\caption{Зависимость мощности работы рассеяния на установившейся стадии ползучести от

напряжения в двойных логарифмических координатах: 

{\sl 1} "--- растяжение образцов, вырезанных вдоль проката; 

{\sl 2} "--- сжатие образцов вдоль проката; 

{\sl 3} "--- сжатие образцов поперек проката; 

{\sl 4} "--- растяжение образцов поперек проката

\label{iya-2.eps}

 }



\smallskip \footnotesize



[Figure~\ref{iya-2.eps}.

Dissipation energy under secondary creep and stresses in double logarithmical coordinates as the dependence curve: 

{\sl 1}---tension  specimens cut out along the direction of the rolling direction of plate; 

{\sl 2}---compression specimens cut out along the direction of the rolling direction of plate; 

{\sl 3}---compression specimens cut out crosswise of the rolling direction of plate; 

{\sl 4}---tension specimens cut out crosswise of the \centerline{rolling direction of plate]}





\end{minipage} \hfill

\begin{minipage}[h]{0.42\textwidth}





\centering

\includegraphics[scale=1.6]{iya-2}



\end{minipage}



\end{figure}



На рис.~\ref{iya-2.eps} представлена зависимость мощности работы рассеяния от напряжения на установившейся стадии ползучести в двойных логарифмических координатах. Диаграмма построена по данным одноосных испытаний, которые обозначены маркерами. Пунктирными линиями приведены зависимости отдельно для растяжения и для сжатия, которые построены на основе степенного закона ползучести с $n=(n_{1}+n_{2})/2$ и соответствующими уже для этого $n$ значениями параметров $B_{A_i}$. 





В табл.~\ref{un:tabl}, \ref{un:tabl2} приведены значения параметров модели ползучести на растяжение и сжатие соответственно, которые получены по данным одноосных испытаний (величина $h$ для $B_{A_i}$ и $B_{\omega_{i}}$ ($i=1,2$) "--- время в~часах). 

Параметры модели ползучести определялись по алгоритму, приведенному в~\cite{iya:Gorev1994}. 

Из данных таблиц и рис.~\ref{iya-2.eps} следует, что для описания поведения материала при данной температуре необходимо использовать модель ползучести, которая описывает различное сопротивление растяжению и~сжатию. 

В~работе~\cite{iya:Tsvelodub2012} представлен вариант модели ортотропной ползучести c~различными свойствами при растяжении и~сжатии. Автором используется степенная форма закона ползучести без учета поврежденности материала. 

Учет различия свойств при растяжении и~сжатии производится путем представления потенциала ползучести в~виде полусуммы потенциалов для растяжения и~для сжатия, а~также добавления слагаемого, которое учитывает знак первого инварианта тензора напряжений. 

Модель~\eqref{iya:sist} является частным случаем модели, рассмотренной в \cite{iya:Tsvelodub2012}, отличие заключается в~учете поврежденности материала и~отсутствии учета ортотропии.









Кроме описанных выше экспериментальных исследований, проведены два испытания на знакопеременный чистый изгиб прямоугольных балок из АБВТ-20 при температуре 750\,\textcelsius. 

Балки вырезались поперек проката плиты и имели следующие геометрические размеры: 

\begin{itemize}

\item[--]балка I: $b = 20.77$~мм, $h = 18.44$~мм; 

\item[--]балка II: $b = 19.60$~мм, $h = 18.40$~мм.

\end{itemize}



Испытания проводились на установке для чистого изгиба балок в~ИГиЛ СО РАН, схема установки приведена в~\cite{iya:Larichkin}. Изменение прогиба балки посредством системы замеров фиксировалось электрическим датчиком линейных перемещений с~последующей записью в~память компьютера. Частота сбора данных: один раз в три секунды. База измерений прогиба составляла 100~мм. 

Программы нагружений изгибающими моментами были следующими: 

\begin{itemize}

\item[--]балка I: $M_{\rm I} = 52.79$~H${}\cdot{}$м, $M_{\rm II} = -52.79$~H${}\cdot{}$м, $M_{\rm III} = -52.79$~H${}\cdot{}$м, $M_{\rm IV} = 52.79$~H${}\cdot{}$м, $M_{\rm V} = 52.79$~H${}\cdot{}$м, $M_{\rm VI} = -52.79$~H${}\cdot{}$м; 

\item[--]балка II: $M_{\rm I} = 36.30$~H${}\cdot{}$м, $M_{\rm II} = -36.30$~H${}\cdot{}$м, $M_{\rm III} = -72.59$~H${}\cdot{}$м, $M_{\rm IV} = 72.59$~H${}\cdot{}$м, $M_{\rm V} = 72.59$~H${}\cdot{}$м.

\end{itemize}







\begin{table}[t!]



\small \centering



\caption{Значения параметров степенного закона ползучести на растяжение\label{un:tabl} }\vspace{-3mm}

[Power law  creep parameters under tension] 

\renewcommand{\tabcolsep}{0.35cm}

\begin{tabular}{p{3.4cm}|c|p{1.8cm}|c|c|p{1.8cm}} \hline



{\footnotesize Направление проката [Rolling direction] }  & \rule{0mm}{12pt} {\footnotesize $n_1$ }  & {\footnotesize  $B_{A_{1}} \cdot 10^{7}$, \par (MPa)$^{-n_1}/$h} & \footnotesize $r_1$& $m_1$& {\footnotesize $B_{\omega_{1}} \cdot 10^{6}$, \par (MPa)$^{-k_1}/$h}\\

            \hline

            \rule{0mm}{12pt}%

           вдоль [along]      &  3.82 &  ~~~~~11.08  &2.72  & 1.34& ~~~~~1.55  \\

           поперек [crosswise]&  3.46 &  ~~~~~79.29  &2.92  & 0.81& ~~~~~3.56 \\

        

        

            \hline

        \end{tabular}



\vspace{2mm}



\small \centering

\caption{Значения параметров степенного закона ползучести  на сжатие\label{un:tabl2} }\vspace{-3mm}

[Power law  creep parameters under compression] 

\renewcommand{\tabcolsep}{0.3cm}

\begin{tabular}{p{3.4cm}|c|p{1.8cm}|c|c|p{1.8cm}} \hline



{\footnotesize Направление проката [Rolling direction] } & \rule{0mm}{12pt}  {\footnotesize $n_2$} & {\footnotesize $B_{A_{2}} \cdot 10^{7}$, \par (MPa)$^{-n_2}/$h}& \footnotesize $r_2$& \footnotesize $m_2$ & {\footnotesize $B_{\omega_{2}} \cdot 10^{7}$, \par (MPa)$^{-n_2}/$h}\\

            \hline

            \rule{0mm}{12pt}%

           вдоль [along]      &  3.86&   ~~~~~4.54&~~~0~~~&~~~0~~~&~~~~~~~~0 \\

           поперек [crosswise]&  4.06& ~~~~~0.92  &~~~0~~~&~~~0~~~&~~~~~~~~0 \\

        

        

            \hline

        \end{tabular}



\end{table}







На рис.~\ref{iya-3},~{\sl a} приведена зависимость прогиба балки I при знакопеременном чистом изгибе. 

На каждом цикле нагружения балка испытывала действие постоянного по модулю момента, после достижения прогиба $\delta_0 \approx 4$~мм производились снятие момента и~дальнейшая выдержка при температуре испытания без нагрузки. После снятия нагрузки датчик прогиба записывал процесс обратной ползучести. 

На рис.~\ref{iya-3},~{\sl b} представлена зависимость модуля прогиба от времени для каждого цикла. 

Начальные точки графиков совмещены в начале координат, поэтому можно видеть уменьшение скорости прогиба балки с~ростом номера цикла нагружения. 

Стоит отметить, что модель ползучести, представленная в~настоящей работе, не учитывает эффект обратной ползучести. Поэтому в дальнейшем~для теоретической обработки результатов использовались лишь кривые ползучести при активном нагружении, а~деформация обратной ползучести не учитывалась.





\begin{figure}[b!]





\footnotesize





\centerline{\begin{tabular}{ccc}

\includegraphics[scale=1.6]{iya-3a} & &

\includegraphics[scale=1.6]{iya-3b} \\

\sl  a & & \sl  b \\

\end{tabular} 

}



\caption{Чистый знакопеременный изгиб балки I из титанового сплава АБВТ-20 при температуре 750\,\textcelsius:

({\sl a}\/)~зависимость абсолютного прогиба от длительности эксперимента; ({\sl b}\/)~зависимость прогиба от времени в каждом цикле 

\label{iya-3}}



\smallskip \footnotesize

[Figure~\ref{iya-3}.

Pure reversed bending of  type I beam made of the titanium ABVT-20 alloy at  750\,\textcelsius: 

({\sl a}\/)~the dependence absolute deflection at duration of experiment;

({\sl b}\/)~the dependence deflection at time in each cycle; the test program:   

$M_{\rm I} = 52.79$~$\rm N \cdot m$,  

$M_{\rm II} = -52.79$~$\rm N \cdot m$, 

\centerline{$M_{\rm III} = -52.79$~$\rm N \cdot m$,  

$M_{\rm IV} = 52.79$~$\rm N \cdot m$, 

$M_{\rm V} = 52.79$~$\rm N \cdot m$,  

$M_{\rm VI} = -52.79$~$\rm N \cdot m$]}



\vspace{-3mm}



\end{figure}



\smallskip





\Section{Сравнение численного расчета с данными испытаний на балках} 

Результаты численного решения~\eqref{iya:eq6}–\eqref{iya:eq7} сравнивались с~данными испытаний знакопеременного чистого изгиба балок. 

На рис.~\ref{iya-4.eps} представлены результаты первых этапов нагружения  балок I ({\sl a}\/) и II ({\sl b}\/). Точки соответствуют экспериментальным зависимостям кривизны балки от времени нагружения $\varkappa=\varkappa(t)$. Черные кривые представляют расчетные значения кривизны, вычисленные по выражению \eqref{iya:eq7}. 

Анализ данной зависимости показывает, что расчетная скорость роста кривизны больше экспериментальной. 



\begin{figure}[b!]



\vspace{3mm}



\footnotesize





\centerline{\begin{tabular}{ccc}

\includegraphics[scale=1.6]{iya-4a} & &

\includegraphics[scale=1.6]{iya-4b} \\

\sl  a & & \sl  b \\

\end{tabular} 

}



\caption{Зависимость кривизны от времени для экспериментальных (точки) и расчетных данных   без учета (черные линии) и с учетом (красные линии) скорости изменения прогиба от длительности эксперимента соответственно для $M=52.79$~H${}\cdot{}$м ({\sl a}\/), $M=36.30$~Н${}\cdot{}$м ({\sl b}\/) }



\label{iya-4.eps}

\smallskip \footnotesize

[Figure~\ref{iya-4.eps}.

Experimental (points) and calculated curvature (solid lines) as the dependence \centerline{deflection  at time for $M=52.79$~N${}\cdot{}$m ({\sl a}\/),\ $M=36.3$~N${}\cdot{}$m ({\sl b}\/)]}

\end{figure}











В связи с неудовлетворенностью таким описанием процесса была проанализирована зависимость изменения скорости ползучести от времени выдержки при температуре. Такая зависимость для скорости прогиба балки приведена в~\cite{iya:Larichkin}. 

Эта зависимость была использована для корректировки параметров модели ползучести. 

В~результате на рис.~\ref{iya-4.eps} красные кривые представляют расчетные значения кривизны балки с~учетом влияния выдержки при температуре. 

Расчет скорости деформаций ползучести производился с~учетом зависимости параметра $B_{A_i}$ от времени по следующей зависимости~\cite{iya:Larichkin}: 

 \begin{equation}

B_{A_i}(t) = 8 \psi(t)^{n_i-1}/(L^{2} M^{n_i-1}),

\end{equation}

где скорость прогиба $\psi(t)=-0.009\cdot \ln(t)+0.0544$; 

$L=100$~мм "--- длина базы измерений; 

$M=52790$~Н${}\cdot{}$м; 

$n_i$ "--- показатели ползучести при растяжении ($i=1$) и~сжатии ($i=2$). 

Данная зависимость получена на основе анализа экспериментальных данных изменения скорости прогиба балки~I, приведенных на рис.~\ref{iya-3},~{\sl b}. 

Это позволило приблизить расчетные данные к~экспериментальным 

(ошибка на каждом этапе нагружения не превосходит~2\,\%).



На рис.~\ref{iya-5.eps} представлена зависимость кривизны балки от времени для знакопеременного чистого изгиба прямоугольной балки II. Изгибающие моменты на каждом этапе были различными, в связи с чем данные этого испытания являются проверочными для рассматриваемой модели ползучести. Точками обозначены экспериментальные данные, сплошными линиями "--- расчетные зависимости с учетом разносопротивляемости, поврежденности и влияния температурной выдержки. Видно, что модель описывает уменьшение скорости прогиба с увеличением времени выдержки при температуре.





\begin{figure}[h!]

\noindent

~~~~\begin{minipage}[h]{0.5\textwidth}



\vspace{7mm}



\caption{Зависимость кривизны от времени в виде экспериментальных (точки) и расчетных данных (сплошные линии) и под действием моментов на каждом этапе: 

$M_{\rm I}=36.3$~Н${}\cdot{}$м,  

$M_{\rm II}=72.59$~Н${}\cdot{}$м,  

$M_{\rm III}=-36.3$~Н${}\cdot{}$м, 

$M_{\rm IV}\hm =-72.59$~Н${}\cdot{}$м,  

$M_{\rm V}=72.59$~Н${}\cdot{}$м} 

\label{iya-5.eps}

 



\smallskip \footnotesize



[Figure~\ref{iya-5.eps}.

Experimental (points) and calculated curvature (solid lines) as the dependence curve by the action of moments at every step: 

$M_{\rm I}=36.3$~$\rm N\cdot m$,  

$M_{\rm II}=72.59$~$\rm N\cdot m$,  

$M_{\rm III}=-36.3$~$\rm N\cdot m$, 

\centerline{$M_{\rm IV}=-72.59$~$\rm N\cdot m$, 

$M_{\rm V}=72.59$~$\rm N\cdot m$]}





\end{minipage} \hfill

\begin{minipage}[h]{0.5\textwidth}





\centering

\includegraphics[scale=1.6]{iya-5}~~~~~ 



\end{minipage}



\end{figure}







\Section[N]{Выводы}

Проведено экспериментальное исследование и математическое моделирование знакопеременного чистого изгиба балок из титанового сплава АБВТ-20 с учетом длительности пребывания материала при температуре 750\,\textcelsius. 

Данные одноосных испытаний при этой температуре позволили установить, что материал является разносопротивляющимся растяжению и сжатию. 

Испытания при знакопеременном изгибе балок показали падение скорости прогиба в зависимости от длительности пребывания при температуре. Такая зависимость была определена на основе данных знакопеременного изгиба прямоугольной балки под действием постоянного по модулю момента сил. Предложена модификация параметров $B_{A_i}$ ($i=1,2$) в модели ползучести при активном нагружении, учитывающая длительность температурной выдержки, что позволило с удовлетворительной погрешностью описать знакопеременный изгиб балок. Одно из возможных применений модификации модели ползучести, учитывающих длительность температурных выдержек, связано, например, с формоизменением изделий в технологических операциях в условиях ползучести материалов, обладающих наблюдаемыми в данном исследовании свойствами.



\AdditionalContent{Конкурирующие интересы}{Мы заявляем, что у нас нет конфликта интересов в авторстве и публикации этой статьи.} 



\AdditionalContent{Авторский вклад и ответственность}{Все  авторы  принимали  участие  в  разработке  концепции  \mbox{статьи}  и  в  написании  рукописи. Мы несем  полную  ответственность  за  предоставление  окончательной рукописи в печать. Окончательная  версия рукописи была нами одобрена.} 



\AdditionalContent{Финансирование}{Работа выполнена при финансовой поддержке РФФИ (проект № 16--08--00713\_а).}

\AdditionalContent{Благодарность}{Авторы благодарны рецензенту за тщательное прочтение статьи и~ценные предложения и комментарии.}



\begin{thebibliography}{10}



\Bibitem{iya:Ribeiro}

\by Ribeiro~F.~C., Marinho~E.~P., Inforzato~D.~J., Costa~P.~R., Batalha~G.~F. 

\paper Creep age forming: a short review of fundaments and applications 

\jour Journal of Achievements in Materials and Manufacturing Engineering

\yr 2010

\vol 43

\issue 1

\pages 353--361

\elink Available at \url{http://jamme.acmsse.h2.pl/papers_vol43_1/43139.pdf} (August 24, 2018)





\Bibitem{iya:Beal}

\by Beal~J.~D., Boyer~R., Sanders~D.

\paper Forming of titanium and titanium Alloys  

\inbook ASM Handbook 

\bookvol 14B

\volinfo Metalworking: Sheet Forming

\yr 2006

\pages 656--669

\publaddr  ASM International







\Bibitem{iya:add:1}

\paper Effect of pre-deformation on Creep age forming of AA2219 plate: Springback, microstructures and mechanical properties

\by Yang~Y., Zhan~L., Ma~Q. et al.

\jour J.~Mater. Process Technology

\vol 229   

\pages 697--702   

\yr 2016

\crossref{10.1016/j.jmatprotec.2015.10.030}

\isi{000367106000072}



\Bibitem{iya:add:2}

\paper Creep-age forming AA2219 plates with different stiffener designs and pre-form age conditions: Experimental and finite element studies

\by Lam~A.~C.~L., Shi~Z., Yang~H. et al.

\jour J.~Mater. Process Technology

\vol 219   

\pages 155--163   

\yr 2015

\crossref{10.1016/j.jmatprotec.2014.12.012}

\isi{000350522000015}





\Bibitem{iya:add:3}

\paper Effect of pre-deformation on creep age forming of 2219 aluminum alloy: Experimental and constitutive modelling

\by Yang~Y., Zhan~L., Shen~R. et al.

\jour Mater. Sci. Eng. A

\vol 683   

\pages 227--235 

\yr 2017 

\crossref{10.1016/j.msea.2016.12.024}

\isi{000393003900027}







\RBibitem{iya:Oleinikov}

\by Олейников~А.~И., Бормотин~К.~С.

\paper Моделирование формообразования крыловых панелей в режиме ползучести с деформационным старением в решениях обратных задач

\jour Ученые записки КнАГТУ

\yr 2015

\issue II-1(22)

\pages 346--365

\elib{https://elibrary.ru/item.asp?id=23608683}

%\misctransl

%\by Oleinikov~A.~I., Bormotin~K.S.

%\paper Modelirovanie formoobrazovaniya krylovyh panelej v rezhime polzuchesti s deformacionnym stareniem v resheniyah obratnyh zadach \rm[Modeling the process of forming wing panels in the creep mode with strain aging in the solution of inverse problems]

%\jour Uchenye zapiski Komsomol'skogo-na-Amure gosudarstvennogo tekhnicheskogo universiteta

%\yr 2015

%\issue II-1(22)

%\pages 346--365

%\lang In Russian



\RBibitem{iya:Korobeynikov}

\by  Коробейников~С.~Н., Олейников~А.~И., Горев~Б.~В., Бормотин~К.~С.

\paper Математическое моделирование процессов ползучести металлических изделий из материалов, имеющих разные свойства при растяжении и сжатии

\jour Вычислительные методы и программирование: Новые вычислительные технологии 

\yr 2008

\vol 9

\issue 1

\pages 346--365

\elib{https://elibrary.ru/item.asp?id=26737529}

%\misctransl

%\by Korobeynikov~S.~N., Oleinikov~A.~I., Gorev~B.~V.,  Bormotin~K.~S. 

%\paper Matematicheskoe modelirovanie processov polzuchesti metallicheskih izdelij iz materialov, imeyushchih raznye svojstva pri rastyazhenii i szhatii \rm[Mathematical simulation of creep processes in metal patterns made of materials with different extension compression properties]

%\jour Vychislitel'nye metody i programmirovanie: Novye vychislitel'nye tekhnologii 

%\yr 2008

%\vol 9

%\issue 1

%\pages 346--365

%\lang In Russian





\RBibitem{iya:Lok2016}

\by Локощенко~А.~М.

\book Ползучесть и длительная прочность металлов

\yr 2016

\publ Физматлит

\publaddr М.

\totalpages 504

%\misctransl

%\by Lokoshenko~A.~M.

%\book Polzuchest' i dlitel'naya prochnost' metallov   \rm [Creep and long-term strength of metals ]

%\yr 1986

%\publ Fizmalit

%\publaddr Moscow

%\lang In Russian

%\totalpages 504







\RBibitem{iya:Kaibishev}

\by Кайбышев~О.~А.

\book Сверхпластичность промышленных сплавов

\yr 1984

\publ Металлургия

\publaddr М.

\totalpages 264

%\misctransl

%\by Kaibishev~O.~A.

%\book Sverhplastichnost' promyshlennyh splavov \rm [The superplasticity of industrial alloys]

%\yr 1984

%\publ Metallurgy

%\publaddr Moscow

%\lang In Russian

%\totalpages 264







\RBibitem{iya:Safiullin}

\by Сафиуллин~Р.~В.

\paper Сверхпластическая формовка и сварка давлением многослойных полых конструкций. Часть II. Опыт ИПСМ РАН

\jour Письма о материалах

\yr 2012

\vol 2

\issue 1 

\pages 36--39

\elib{https://elibrary.ru/item.asp?id=17728833}

%\misctransl

%\by Safiullin~R.~V.

%\paper Sverhplasticheskaya formovka i svarka davleniem mnogoslojnyh polyh konstrukcij. Chast' II. Opyt IPSM RAN \rm[Superplastic forming and pressure welding of multilayer

%hollow structures Part II. Experience of IMSP RAS]

%\jour Pis'ma o materialah

%\yr 2012

%\issue 2

%\pages 36--39

%\lang In Russian



\RBibitem{iya:Radchenko}

\by Радченко~В.~П., Саушкин~М.~Н.

\book Ползучесть и релаксация  остаточных напряжений в упрочненных конструкциях

\yr 2005

\publ Машиностроение-1

\publaddr М.

\totalpages 226

%\misctransl

%\by Radchenko~V.~P., Saushkin.~M.~N.

%\book Polzuchest' i relaksaciya  ostatochnyh napryazhenij v uprochnennyh konstrukciyah \rm [Creep and relaxation of residual stresses in hardened structures]

%\yr 2005

%\publ Metallurgy-1

%\publaddr Moscow

%\lang In Russian

%\totalpages 226



\RBibitem{iya:Loko2012}

\by  Локощенко~А.~М., Агахи~К.~А., Фомин~Л.~В. 

\paper Изгиб балки при ползучести с учетом поврежденности и разносопротивляемости материала 

\jour Машиностроение и инженерное образование

\yr 2012

\issue 3(32)

\pages 29--35

\elib{https://elibrary.ru/item.asp?id=21118278}

%\misctransl

%\by Lokoshchenko~A.~M., Agakhi~K.~A., Fomin~L.~V.

%\paper Izgib balki pri polzuchesti s uchetom povrejdaemosti i raznosoprotivlyaemosti materiala \rm [Beam bending under creep considering the damage and multimodulus behavior of the material]

%\jour Mechanic engineering and engineering education 

%\yr 2012

%\issue 3(32)

%\pages 29--35

%\lang In Russian



\RBibitem{iya:Larichkin}

\by  Колодезев~В.~Е., Горев~Б.~В., Ларичкин~A.~Ю., Шевцова~Л.~И.

\paper Чистый изгиб балки из сплава АБВТ-20 в режиме ползучести при знакопеременном нагружении

\jour Технология машиностроения

\yr 2017

\issue 2(176)

\pages 11--16

\elib{https://elibrary.ru/item.asp?id=29715192}

%\misctransl

%\by Kolodezev~V.~E., Gorev~B.~V., Larichkin~A.~Yu.,  Shevtsova~L.~I. 

%\paper Chistii izgib balki iz splava ABVT-20 v rejime polzuchesti pri znakoperemennom nagrujenii \rm[Pure bending of beams from an alloy ABVT-20 in creep mode at alternating loading]

%\jour Tekhnologiya Mashinostroeniya

%\yr 2017

%\issue 2(176)

%\pages 11--16

%\lang In Russian



\RBibitem{iya:Sosnin}

\by Соснин~О.~В., Горев~Б.~В., Никитенко~А.~Ф.  

\book Энергетический вариант теории ползучести

\yr 1986

\publ ИГиЛ СО АН СССР

\publaddr Новосибирск

\totalpages 95

%\misctransl

%\by Sosnin ~O.~V., Gorev ~B.~V., Nikitenko.~A.~F. 

%\book Energeticheskii  Variant Teorii Polzuchesti   \rm [Energy Variant of the Creep Theory]

%\yr 1986

%\publ Lavrentyev Institute of Hydrodynamics SO AN SSSR

%\publaddr Novosibirsk

%\lang In Russian

%\totalpages 95



















\RBibitem{iya:Nikitenko}

\by Никитенко~А.~Ф.

\book Ползучесть и длительная прочность металлических материалов

\yr 1997

\publ Новосиб. гос. архит.-строит. ун-т

\publaddr Новосибирск

\totalpages 278

%\misctransl

%\by Nikitenko~A.~F.

%\book Polzuchest’ i dlitel’naya prochnost’ metallicheskikh materialov  \rm [Creep and Long-term Strength of Metal Materials]

%\yr 1997

%\publ Novosibirsk State University of Architecture and Civil Engineering

%\publaddr Novosibirsk

%\lang In Russian

%\totalpages 278







\RBibitem{iya:Gorev2014}

\by Горев~Б.~В., Любашевская~И.~В., Панамарев~В.~А., Иявойнен~С.~В.

\paper Описание процесса  ползучести и разрушения  современных конструкционных материалов с~использованием кинетических уравнений в энергетической форме 

\jour ПМТФ

\yr 2014

\vol 55

\issue 6

\pages 132--144

%\misctransl

%\crossref{10.1134/S0021894414060145}

%\by Gorev~B.~V., Lubashevskaya~I.~V., Panamarev~V.~A., Iyavoynen~S.V.

%\paper Opisanie processa  polzuchesti i razrusheniya  sovremennyh konstrukcionnyh materialov s ispol'zovaniem kineticheskih uravnenij v energeticheskoj forme \rm[Description of creep and fracture of modern construction materials using kinetic equations in energy form]

%\jour Journal of Applied Mechanics and Technical Physics

%\yr 2014

%\vol 55

%\issue 6

%\pages 132--144

%\lang In Russian





\RBibitem{iya:Gorev1973}

\by  Горев~Б.~В.

\paper К расчету на неустановившуюся ползучесть изгибаемого бруса из материала с разными характеристиками на растяжение и сжатие 

\inbook Динамика сплошной среды: сб. науч. тр.

\issueinfo Вып. 14

\publ АН СССР. Сиб. отд-ние. Ин-т гидродинамики

\publaddr Новосибирск

\yr 1973

\pages 44--51

%\misctransl

%\by Gorev~B.~V.

%\paper  K raschetu na neustanovivshuyusya polzuchest’ izgibaemogo brusa iz materiala s raznymi kharakteristikami na rastyazhenie i szhatie \rm [To calculation for transient creep beam bending of material with different characteristics in tension and compression]

%\jour Continuum Dynamics

%\yr 1973

%\issue 14

%\pages 44--51

%\lang In Russian







\RBibitem{iya:Kuznetsov}

\by  Кузнецов~Е.~Б., Леонов~С.~С. 

\paper Чистый изгиб балки из разномодульного материала в условиях ползучести

\jour Вестн. ЮУрГУ. Сер. Матем. моделирование и программирование

\yr 2013

\vol 6

\issue 4

\pages 26--38

\mathnet{http://mi.mathnet.ru/vyuru101}





\RBibitem{iya:Sosnin1970}

\by  Соснин~О.~В. 

\paper О ползучести материалов с разными характеристиками на растяжение и сжатие 

\jour ПМТФ

\yr 1970

\issue 5

\pages 136--139

%\misctransl

%\by Sosnin~O.~V.

%\paper O polzuchesti materialov s raznimi charachteristikami na rastyajenie i sjatie \rm [On creep of materials with different tension and compression properties]

%\jour Journal of Applied Mechanics and Technical Physics

%\yr 1970

%\issue 5

%\pages 136--139

%\lang In Russian











\Bibitem{iya:Runge}

\by Merson~R.~H.

\paper An operational method for the study of integration processes

\inbook Proc. Symp. Data Processing, Weapons Res. Establ. Salisbury

\yr 1957

\pages 110–125

\publaddr Salisbury

\moreref

\by \linebreak Pospelov~V.~V.

\preprint  Kutta-Merson method

\preprintinfo Encyclopedia of Mathematics

\yr 2014

\elink Available at \url{http://www.encyclopediaofmath.org/index.php?title=Kutta-Merson_method&oldid=33669} (August~24, 2018)







\Bibitem{iya:Simpson}

\by Pao~Y.~C.

\book Engineering analysis. Interactive methods and programs with FORTRAN, QuickBASIC, MATLAB, and Mathematica

\zmath{0908.65001}

\publaddr Boca Raton, FL

\publ CRC Press

\totalpages 360 

\yr 1998





\RBibitem{iya:Rabotnov}

\by Работнов~Ю.~Н.

\book Ползучесть элементов конструкций

\yr 2014

\publ Наука

\publaddr М.

\totalpages 752

%\misctransl

%\by Rabotnov~Yu.~N.

%\book Polzuchest' elemontov construczii \rm [ Creep of construction elements ]

%\yr 2014

%\publ Nauka

%\publaddr Moscow

%\lang In Russian

%\totalpages 752





\RBibitem{iya:Gorev1994}

\by  Горев~Б.~В., Клопотов~И.~Д. 

\paper К описанию процесса ползучести и длительной прочности по уравнениям с одним скалярным параметром повреждаемости

\jour ПМТФ

\yr 1994

\issue 5

\pages 92--102

%\misctransl

%\by Gorev~B.~V., Klopotov~I.~D.

%\paper K opisaniju processa polzuchesti i dlitel’noj prochnosti po uravnenijam s odnim skaljarnym parametrom povrezhdaemosti \rm [Description of the creep process and long-term strength by equations with one scalar damage parameter]

%\jour Journal of Applied Mechanics and Technical Physics

%\yr 1994

%\issue 5

%\pages 92--102

%\lang In Russian



\RBibitem{iya:Tsvelodub2012}

\by  Цвелодуб~И.~Ю. 

\paper К построению определяющих уравнений ползучести ортотропных материалов с различными свойствами при растяжении и сжатии

\jour ПМТФ

\yr 2012

\issue 6

\pages 98--101

\elib{https://elibrary.ru/item.asp?id=18194293}

%\misctransl

%\by Tsvelodub~I~Y.

%\paper K postroeniju opredeljajushih uravnenij polzuchesti ortotropnych materialov s razlichnymi svojstvami pri rastjazhenii i szhatii \rm [Construction of constitutive equations of creep in orthotropic materials with different properties under tension and compression]

%\jour Journal of Applied Mechanics and Technical Physics

%\yr 2012

%\issue 6

%\pages 98--101

%\lang In Russian



\end{thebibliography}



\renewcommand{\baselinestretch}{0.89}



\newpage







\makeentitle



\vspace{-5mm}



\AdditionalContent{Competing interests}{We declare that we have no conflicts of interest in the authorship and publication of this article.} 



\AdditionalContent{Authors' contributions and responsibilities}{Each author has participated in the article concept development and in the manuscript writing. The authors are absolutely responsible for submitting the final manuscript in print. Each author has approved the final version of manuscript.} 



\AdditionalContent{Funding}{This work was supported by the Russian Foundation for Basic Research (project no. 16--08--00713\_a).}

\AdditionalContent{Acknowledgments}{The authors are grateful to the  referee for  careful reading  of  the paper and valuable  suggestions and comments.}







\begin{thebibliography}{10}



\Bibitem{iya:Ribeiro:en}

\by Ribeiro~F.~C., Marinho~E.~P., Inforzato~D.~J., Costa~P.~R., Batalha~G.~F. 

\paper Creep age forming: a short review of fundaments and applications 

\jour Journal of Achievements in Materials and Manufacturing Engineering

\yr 2010

\vol 43

\issue 1

\pages 353--361

\elink Available at \url{http://jamme.acmsse.h2.pl/papers_vol43_1/43139.pdf} (August 24, 2018)





\Bibitem{iya:Beal:en}

\by Beal~J.~D., Boyer~R., Sanders~D.

\paper Forming of titanium and titanium Alloys  

\inbook ASM Handbook 

\bookvol 14B

\volinfo Metalworking: Sheet Forming

\yr 2006

\pages 656--669

\publaddr  ASM International







\Bibitem{iya:add:1:en}

\paper Effect of pre-deformation on Creep age forming of AA2219 plate: Springback, microstructures and mechanical properties

\by Yang~Y., Zhan~L., Ma~Q. et al.

\jour J.~Mater. Process Technology

\vol 229   

\pages 697--702   

\yr 2016

\crossref{10.1016/j.jmatprotec.2015.10.030}

\isi{000367106000072}



\Bibitem{iya:add:2:en}

\paper Creep-age forming AA2219 plates with different stiffener designs and pre-form age conditions: Experimental and finite element studies

\by Lam~A.~C.~L., Shi~Z., Yang~H. et al.

\jour J.~Mater. Process Technology

\vol 219   

\pages 155--163   

\yr 2015

\crossref{10.1016/j.jmatprotec.2014.12.012}

\isi{000350522000015}





\Bibitem{iya:add:3:en}

\paper Effect of pre-deformation on creep age forming of 2219 aluminum alloy: Experimental and constitutive modelling

\by Yang~Y., Zhan~L., Shen~R. et al.

\jour Mater. Sci. Eng. A

\vol 683   

\pages 227--235 

\yr 2017 

\crossref{10.1016/j.msea.2016.12.024}

\isi{000393003900027}







\Bibitem{iya:Oleinikov:en}

\lang In Russian

\by Oleinikov~A.~I., Bormotin~K.S.

\paper Modeling the process of forming wing panels in the creep mode with strain aging in the solution of inverse problems

\jour Uchenye Zapiski Komsomol'skogo-na-Amure Gosudarstvennogo Tekhnicheskogo Universiteta

\yr 2015

\issue II-1(22)

\pages 346--365





\Bibitem{iya:Korobeynikov:en}

\lang In Russian

\by Korobeynikov~S.~N., Oleinikov~A.~I., Gorev~B.~V.,  Bormotin~K.~S. 

\paper Mathematical simulation of creep processes in metal patterns made of materials with different extension compression properties

\jour Vychislitel'nye Metody i Programmirovanie: Novye Vychislitel'nye Tekhnologii

\yr 2008

\vol 9

\issue 1

\pages 346--365







\Bibitem{iya:Lok2016:en}

\by Lokoshchenko~A.~M.

\book Polzuchest' i dlitel'naia prochnost' metallov \rm [Creep and Long-Term Strength of Metals]

\yr 2016

\publ Fizmatlit

\publaddr Moscow

\totalpages 504

\lang In Russian









\Bibitem{iya:Kaibishev:en}

\lang In Russian

\by Kaibyshev~O.~A.

\book Sverkhplastichnost' promyshlennykh splavov  \rm [The Superplasticity of Industrial Alloys]

\yr 1984

\publ Metallurgiia

\publaddr Moscow

\totalpages 264







\Bibitem{iya:Safiullin:en}

\lang In Russian

\by Safiullin~R.~V.

\paper Superplastic forming and pressure welding of multilayer hollow structures. Part II. Experience of IMSP RAS

\jour Pis'ma o Materialakh

\yr 2012

\vol 2

\issue 1 

\pages 36--39





\Bibitem{iya:Radchenko:en}

\lang In Russian

\by Radchenko~V.~P., Saushkin~M.~N.

\book Polzuchest' i relaksatsiia  ostatochnykh napriazhenii v uprochnennykh konstruktsiiakh \rm [Creep and Relaxation of Residual Stresses in Hardened Structures]

\yr 2005

\publ Mashinostroenie-1

\publaddr Moscow

\totalpages 226





\Bibitem{iya:Loko2012:en}

\lang In Russian

\by Lokoshchenko~A.~M., Agakhi~K.~A., Fomin~L.~V.

\paper Beam bending under creep considering the damage and multimodulus behavior of the material

\jour Mechanic Engineering and Engineering Education 

\yr 2012

\issue 3(32)

\pages 29--35





\Bibitem{iya:Larichkin:en}

\lang In Russian

\by  Kolodezev~V.~E., Gorev~B.~V., Larichkin~A.~Iu., Shevtsova~L.~I.

\paper Pure bending of beams from an alloy ABVT-20 in creep mode at alternating loading

\jour Tekhnologiia Mashinostroeniia

\yr 2017

\issue 2(176)

\pages 11--16





\Bibitem{iya:Sosnin:en}

\lang In Russian

\by Sosnin~O.~V., Gorev~B.~V., Nikitenko~A.~F.  

\book Energeticheskii variant teorii polzuchesti \rm [Energy Variant of the Creep Theory]

\yr 1986

\publ Lavrentiev Institute of Hydrodynamics

\publaddr Novosibirsk

\totalpages 95



















\Bibitem{iya:Nikitenko:en}

\lang In Russian

\by Nikitenko~A.~F.

\book Polzuchest' i dlitel'naia prochnost' metallicheskikh materialov \rm [Creep and Long-term Strength of Metal Materials]

\yr 1997

\publ Novosibirsk State University of Architecture and Civil Engineering

\publaddr Novosibirsk

\totalpages 278









\Bibitem{iya:Gorev2014:en}

\by Gorev~B.~V., Lubashevskaya~I.~V., Panamarev~V.~A., Iyavoynen~S.V.

\jour J. Appl. Mech. Tech. Phys. 

\yr 2014

\crossref{10.1134/S0021894414060145}

\yr 2014

\vol 55

\issue 6

\pages 1020–1030 

\paper Description of creep and fracture of modern construction materials using kinetic equations in energy form







\Bibitem{iya:Gorev1973:en}

\lang In Russian

\by  Gorev~B.~V.

\paper To calculation for transient creep beam bending of material with different characteristics in tension and compression

\inbook Dinamika sploshnoi sredy \rm [Continuum Dynamics]

\issueinfo Issue 14

\publ Lavrentiev Institute of Hydrodynamics

\publaddr Novosibirsk

\yr 1973

\pages 44--51









\Bibitem{iya:Kuznetsov:en}

\lang In Russian

\by Kuznetsov~E.~B., Leonov~S.~S.

\paper Pure Bending for the Multimodulus Material Beam at Creep Conditions

\jour Vestnik YuUrGU. Ser. Mat. Model. Progr.

\yr 2013

\vol 6

\issue 4

\pages 26--38





\Bibitem{iya:Sosnin1970:en}

\yr 1970

\vol 11

\issue 5

\pages 832–835 

\paper Creep in materials with different tension and compression behavior

\by Sosnin~O.~V.

\jour  J.~Appl. Mech. Tech. Phys.

\crossref{10.1007/BF00851914}











\Bibitem{iya:Runge:en}

\by Merson~R.~H.

\paper An operational method for the study of integration processes

\inbook Proc. Symp. Data Processing, Weapons Res. Establ. Salisbury

\yr 1957

\pages 110–125

\publaddr Salisbury

\moreref

\by \linebreak Pospelov~V.~V.

\preprint Kutta-Merson method

\preprintinfo Encyclopedia of Mathematics

\yr 2014

\elink Available at \url{http://www.encyclopediaofmath.org/index.php?title=Kutta-Merson_method&oldid=33669} (August~24, 2018)







\Bibitem{iya:Simpson:en}

\by Pao~Y.~C.

\book Engineering analysis. Interactive methods and programs with FORTRAN, QuickBASIC, MATLAB, and Mathematica

\zmath{0908.65001}

\publaddr Boca Raton, FL

\publ CRC Press

\totalpages 360 

\yr 1998





\Bibitem{iya:Rabotnov:en}

\lang In Russian

\by Rabotnov~Yu.~N.

\book Polzuchest' elementov konstruktsii \rm [Creep of Construction Elements]

\yr 2014

\publ Nauka

\publaddr Moscow

\totalpages 752





\Bibitem{iya:Gorev1994:en}

\by Gorev~B.~V., Klopotov~I.~D.

\jour J.~Appl. Mech. Tech. Phys.

\yr 1994

\crossref{10.1007/BF02369552}

\vol 35

\issue 5

\pages 726–734 

\paper Description of the creep process and long-term strength by equations with one scalar damage parameter





\Bibitem{iya:Tsvelodub2012:en}

\by Tsvelodub~I.~Yu.

\jour J.~Appl. Mech. Tech. Phys 

\yr 2012

\crossref{10.1134/S0021894412060119}

\vol 53

\issue 6

\pages 888–890 

\paper Construction of constitutive equations of creep in orthotropic materials with different properties under tension and compression





\end{thebibliography}







\renewcommand{\baselinestretch}{0.9}



%\end{document}

